\section{Space, alternation and interaction}

\begin{result}
\textbf{Containment map.}
The proved machine-class relationships include
\[
\cc{NL}\subseteq\cc{P}\subseteq\cc{NP}\subseteq\cc{PH}
\subseteq\cc{PSPACE}\subseteq\cc{EXP},
\qquad \cc{AP}=\cc{NPSPACE}=\cc{PSPACE}.
\]
The equalities and the NL inclusion receive entries below. The links here supply
the remaining arrows, including PH containment in polynomial space and the
deterministic exponential-time simulation of polynomial space. Inclusion arrows do not assert strictness;
the separate $\cc{P}\subsetneq\cc{EXP}$ theorem appears in the preceding section.
\citep{AroraBarak09,Sipser2013}
\decl{Complexity/Computability/Complexity/Classes/Ladner.lean}{103}{Cslib.Complexity.P_subset_NP}
\decl{Complexity/Computability/Complexity/Classes/PH.lean}{295}{Cslib.Complexity.NP_subset_PH}
\decl{Complexity/Computability/Complexity/Classes/Inclusions.lean}{572}{Cslib.Complexity.PH_subset_PSPACE}
\decl{Complexity/Computability/Complexity/Classes/Inclusions.lean}{260}{Cslib.Complexity.PSPACE_subset_EXP}
\end{result}

\begin{result}
\subsection{Nondeterministic logspace lies in polynomial time}
\[
\cc{NL}\subseteq\cc{P}.
\]
NL is defined by native nondeterministic machines with logarithmic space bounded
on every computation prefix and acceptance by an accepting path. A polynomial-time
decider exists even though the nondeterministic machine's other branches need not
halt. This entry does not assert the still-outstanding native $\cc{NL}=\cc{coNL}$
theorem.\citep[Section 8.5]{Sipser2013}
\decl{Complexity/Computability/Complexity/Classes/LogSpace/PolyTime.lean}{44}{Cslib.Complexity.NL_subset_P}
\end{result}

\begin{result}
\subsection{Undirected connectivity by short walks in Reingold's graph}
For every simple graph $G$ on $\{0,\dots,n-1\}$ and all vertices $s,t$, $t$ is reachable from $s$
if and only if some walk of length $T(n)=4(2\lceil\log_2 n\rceil+65\,\ell(n))$ in the
$257^8$-regular graph $G_{\ell(n)}$, starting at a fixed vertex over the slot $(s,s)$, ends over a
slot of $t$, where $\ell(n)=16\lceil\log_2 n\rceil+524$. Here $G_0$ is the slot graph of $G$ (each
vertex replaced by a cycle of $n$ slots, slot $(v,w)$ joined to $(w,v)$ when $v\sim w$), padded
with self-loops to $257^8$ labels, and $G_{i+1}=((G_i\mathbin{\textcircled{\scriptsize z}}H)^2)$ with labels renamed, for
the arithmetic zig-zag base $H$ over $\mathbb{Z}/257$ with $\lambda(H)\le 1/5$. The generic
statement assumes $\lambda(H)\le1/5$, at least half self-loops at each vertex,
$2|V|^2|B|^2 5^\ell\le 5\cdot 6^\ell$ and $|V||B|^\ell 4^T\le 5^T$. Each step turns a spectral gap
$\gamma$ into at least $\min(6\gamma/5,1/5)$ (Reingold takes the eighth power of $G\mathbin{\textcircled{\scriptsize z}}H$
and obtains $\lambda\mapsto\max(\lambda^2,1/2)$; the constant differs, the $O(\log n)$ depth does
not), a connected component with half self-loops has gap at least $1/(2N^2|B|^2)$, and walks of
length $t$ connect every pair once $\lambda^t(N-1)<1$. This is the combinatorial part of the proof
of $\mathrm{USTCON}\in\cc{L}$; log-space evaluation of the rotation map and the class-level theorem
are not yet formalized.\citep[Definition 3.1, Lemmas 3.2--3.3, Theorem 4.1]{Reingold2008}\citep[Section 21.4]{AroraBarak09}
\declat{cb631cc99441574d663dc693899a59a5b1e906f3}{Complexity/Algorithms/Reachability/Reingold.lean}{223}{Cslib.Reachability.reachable_iff_exists_reingoldGraph}
\declat{cb631cc99441574d663dc693899a59a5b1e906f3}{Complexity/Foundations/Combinatorics/Expander/Reingold.lean}{428}{Cslib.RotationMap.reachable_iff_exists_root_neighbor_pow_reingold}
\declat{cb631cc99441574d663dc693899a59a5b1e906f3}{Complexity/Foundations/Combinatorics/Expander/Reingold.lean}{318}{Cslib.RotationMap.min_pow_le_one_sub_spectralExpansion_reingold}
\declat{cb631cc99441574d663dc693899a59a5b1e906f3}{Complexity/Foundations/Combinatorics/Expander/Connected.lean}{135}{Cslib.RotationMap.spectralExpansion_le_of_preconnected}
\declat{cb631cc99441574d663dc693899a59a5b1e906f3}{Complexity/Foundations/Combinatorics/Expander/Connected.lean}{182}{Cslib.RotationMap.exists_neighbor_pow_eq}
\end{result}


\begin{result}
\subsection{Savitch's theorem for polynomial space}
\[
\cc{NPSPACE}=\cc{PSPACE}.
\]
Both classes are machine-defined. Nondeterministic acceptance is by an accepting
path, with every path respecting the space bound; the deterministic decider is
total. This is the equality of polynomial-space classes, not a statement of the
fully parameterized $\cc{NSPACE}(s)\subseteq\cc{DSPACE}(s^2)$ theorem.
\citep[Corollary 8.6]{Sipser2013}
\decl{Complexity/Computability/Complexity/Classes/Space/Savitch.lean}{54}{Cslib.Complexity.NPSPACE_eq_PSPACE}
\end{result}

\begin{result}
\subsection{TQBF is PSPACE-complete}
The true closed quantified Boolean formulas are PSPACE-complete under
polynomial-time many-one reductions. Membership uses an explicit-stack evaluator;
hardness includes a certified polynomial-time construction of the quantified
reachability formula. Formulas with free variables are excluded from TQBF.
\citep[Theorem 8.9]{Sipser2013}
\decl{Complexity/Computability/Complexity/Classes/Space/TQBF.lean}{86}{Cslib.Complexity.TQBF_pspaceComplete}
\end{result}

\begin{result}
\subsection{Alternating polynomial time equals polynomial space}
\[
\cc{AP}=\cc{PSPACE}.
\]
AP uses the library's native alternating machines and a polynomial time bound.
The equality is with the same machine-defined PSPACE used above; it is not merely
an equality of abstract game predicates.\citep{ChandraKozenStockmeyer1981}
\decl{Complexity/Computability/Complexity/Classes/Alternation/Savitch.lean}{385}{Cslib.Complexity.AP_eq_PSPACE}
\end{result}

\begin{result}
\subsection{Interactive proofs lie in polynomial space}
\[
\cc{IP}\subseteq\cc{PSPACE}.
\]
IP uses a private-coin verifier realized by one fair-coin oracle machine and a
polynomial clock valid against every stateful prover. Provers may be randomized
and computationally unbounded; completeness is at least $2/3$ and soundness at
most $1/3$. The converse $\cc{PSPACE}\subseteq\cc{IP}$ is not yet formalized here.
\citep[Chapter 8]{AroraBarak09}
\decl{Complexity/Computability/Complexity/Classes/IP/PSPACE.lean}{620}{Cslib.Complexity.IP_subset_PSPACE}
\end{result}

\begin{result}
\subsection{The algebraic protocol for quantified Boolean formulas}
Let $\varphi$ be a closed quantified Boolean formula with $|\varphi|$ syntax nodes, and let
$E$ be its arithmetization over a commutative ring, with Shen's linearization operator
applied to every variable of a list $xs$ after every quantifier. Then $E(0)=1$ if $\varphi$
is true and $E(0)=0$ if it is false. If $xs$ contains the variables of $\varphi$, every
syntactic degree bound of $E$ is at most $|\varphi|$, and the round degree (the sum of these
bounds over the operator nodes) is at most $(2|xs|+1)|\varphi|^2$. A Shamir--Shen-style
verifier checks the claim $E(0)=1$ by recursion on $E$: at a binary connective the prover
sends the two values, and at $\exists x$, $\forall x$ or $L_x$ it sends the univariate
restriction, truncated to the degree bound, and the verifier draws a random point. The
honest prover is accepted with probability $1$ on every true claim. Over an integral
domain, if the sampler gives each element probability at most $\varepsilon$, every
stateful, adaptive, randomized prover is accepted on a false claim with probability at
most (round degree)$\cdot\varepsilon$. For TQBF with uniform samples from a finite field
$F$, false closed formulas are accepted with probability at most $3|\varphi|^3/|F|$. The
verifier is an abstract typed oracle program with exact sampling. Its realization in the
IP machine model, and hence $\cc{PSPACE}\subseteq\cc{IP}$, is not yet formalized.
\citep{LundFortnowKarloffNisan1992,Shamir1992,Shen1992}\citep[Chapter 8]{AroraBarak09}
\declat{b18f43891513f65a64c5417d388abaaee4548ef1}{Complexity/Logics/Propositional/QBF/Arithmetization.lean}{485}{Cslib.Logic.PL.QBF.eval_arithmetize}
\declat{b18f43891513f65a64c5417d388abaaee4548ef1}{Complexity/Logics/Propositional/QBF/Arithmetization.lean}{491}{Cslib.Logic.PL.QBF.eval_arithmetize_zero}
\declat{b18f43891513f65a64c5417d388abaaee4548ef1}{Complexity/Logics/Propositional/QBF/Arithmetization.lean}{523}{Cslib.Logic.PL.QBF.deg_arithmetize_le}
\declat{b18f43891513f65a64c5417d388abaaee4548ef1}{Complexity/Logics/Propositional/QBF/Arithmetization.lean}{590}{Cslib.Logic.PL.QBF.roundDegree_arithmetize_le}
\declat{b18f43891513f65a64c5417d388abaaee4548ef1}{Complexity/Computability/Complexity/Classes/IP/TQBF/Protocol.lean}{338}{Cslib.Complexity.QBFProtocol.acceptProb_honest_verify}
\declat{b18f43891513f65a64c5417d388abaaee4548ef1}{Complexity/Computability/Complexity/Classes/IP/TQBF/Protocol.lean}{347}{Cslib.Complexity.QBFProtocol.acceptProb_verify_le}
\declat{b18f43891513f65a64c5417d388abaaee4548ef1}{Complexity/Computability/Complexity/Classes/IP/TQBF/Protocol.lean}{464}{Cslib.Complexity.QBFProtocol.acceptProb_honest_of_mem_TQBF}
\declat{b18f43891513f65a64c5417d388abaaee4548ef1}{Complexity/Computability/Complexity/Classes/IP/TQBF/Protocol.lean}{471}{Cslib.Complexity.QBFProtocol.acceptProb_le_of_notMem_TQBF}
\end{result}


\begin{result}
\subsection{Graph non-isomorphism has a private-coin interactive proof}
Graph non-isomorphism belongs to IP. The strict polynomial-time fair-coin verifier
has perfect completeness and soundness at most $(17/32)^2<1/3$ after two rounds.
It samples relabelings by sorting random vertex keys; this avoids assuming exact
uniform sampling of a permutation from a fixed number of fair bits. The result is
for private coins and does not establish the public-coin AM protocol.
\citep{GoldreichMicaliWigderson1991}
\decl{Complexity/Computability/Complexity/Classes/IP/GraphNonIsomorphism.lean}{563}{Cslib.Complexity.GraphNonIsomorphism.isInteractiveProof_verifier}
\decl{Complexity/Computability/Complexity/Classes/IP/GraphNonIsomorphism.lean}{603}{Cslib.Complexity.GraphNonIsomorphism.mem_IP}
\end{result}


\begin{result}
\subsection{Cook--Mertz Tree Evaluation in small space}
A valid instance is a perfect binary tree with every leaf at depth $H$, with $\ell$-bit values at
its leaves and a gate
$\{0,1\}^{\ell}\times\{0,1\}^{\ell}\to\{0,1\}^{\ell}$ at each internal node.
The injective encoding gives $\ell,H$ in unary, followed by gate truth tables
and leaf values in preorder. Write $n$ for its bit length and $b(k)$ for the
length of the library's binary encoding of $k$.
One fixed deterministic multi-tape Turing machine computes the root value on
every encoded instance, using work space
\[
O\bigl((\ell+H+1)(b(\ell)+1)+b(n+1)\bigr),
\]
that is, $O((\ell+H+1)\log(\ell+2)+\log(n+2))$.
The theorem allows zero value length and zero height. Input addressing costs
are included, the output is streamed, and a running-time bound exists without
an asserted asymptotic rate.

Running once over every nonzero field element, the algebraic procedure preserves
the other two registers while adding the
scaled tree value to its target register over any finite field with
$2\ell<|K|-1$. The concrete machine chooses a suitable prime, reads the encoded
gates, and satisfies a proved state-size invariant. This formalizes the
Cook--Mertz construction with Goldreich's value-length parametrization and
exposition; the Lean development is by Christian Reitwiessner with disclosed
AI assistance.\citep{CookMertz2024,Goldreich2024TreeEval}
Correctness and resource bounds apply to valid encoded instances; malformed-input
rejection is not proved. A separate $O(\log n\log\log n)$ corollary and the refined
$O(\ell+H\log(\ell+2))$ parameter bound are not asserted here.
\declat{aa72232abd8858a2c4103eb6aebc039e3b02f52d}{Complexity/Computability/TreeEval/Space.lean}{39}{Cslib.TreeEval.computableTidilyInSpace_eval}
\declat{aa72232abd8858a2c4103eb6aebc039e3b02f52d}{Complexity/Computability/TreeEval/Instance.lean}{74}{Cslib.TreeEval.encode_injective}
\declat{aa72232abd8858a2c4103eb6aebc039e3b02f52d}{Complexity/Computability/TreeEval/CookMertz.lean}{157}{Cslib.TreeEval.proc_eq}
\end{result}
